\documentclass{article}
\usepackage{ctex}
\begin{document}


%高斯单位制当中，力的单位是达因，距离的单位是厘米，电荷量的单位是statC请将 statC 用国际单位制表示出来。

在高斯单位制中

$$1\rm{dyn}=\frac{(1\rm{statC})^2}{(1\rm{cm})^2}$$

$$1\rm{dyn}=1\times10^{-5}\rm{N}\qquad1\rm{cm}=1\times10^{-2}\rm{m}$$

$$(1\rm{statC})^2=1\times10^{-5}\rm{N}\times(1\times10^-2\rm{m})^2=10^9\rm{kg}\cdot \rm{m}^3\cdot\rm{s}^{-2}$$

$$1\rm{statC}=10^{-\frac{9}{2}}\rm{kg}^{\frac{1}{2}}\cdot \rm{m}^{\frac{3}{2}}\cdot\rm{s}^{-1}$$



%1.1-3 为了得到一库仑电荷量大小的概念,试计算两个都是一库仑的点电荷在真空中相距 $1\rm{m}$ 时的相互作用力和相距 $1 \rm{km}$ 时的相互作用力.

1.1-3

库仑力的大小为：

$$F=k\frac{q_1q_2}{r^2}$$

$$k=8.99\times10^9\rm{N}\cdotp \rm{m}^2/ \rm{C}^2,\qquad q_1=q_2=1\rm{C}$$

$$r=1\rm{m},F=8.99\times10^9\rm{N}$$

$$r=1\rm{km},F=8.99\times10^3\rm{N}$$

\newpage
%1.1-4 氢原子由一个质子(即氢原子核)和一个电子组成.根据经典模型,在正常状态下,电子绕核做圆周运动,轨道半径是$5.29x10^{-11}m$.
%已知质子质量$m_p=1.67x10^{-27}\rm{kg}$,电子质量$m_e=9.11x10^{-31}\rm{kg}$,电荷分别为$\pm e=$ $\pm1.60\times10^{-19}$ C,万有引力常量$G=6.67\times10^{-11}\rm{N}\cdot{m}^2\cdot\rm{kg}^2$.

%(1)求电子所受的库仑力和引力；

%(2)库仑力是万有引力的多少倍？

%(3)求电子的速度.

1.1-4

(1)

库仑力的大小为：

$$F_e=k\frac{q_1q_2}{r^2}$$
$$k=8.99\times10^9\rm{N}\cdotp \rm{m}^2/ \rm{C}^2,$$
$$ q_1=q_2=1.60\times10^{-19}\rm{C} $$
$$  r=5.29\times10^{-11}\rm{m}$$
$$F_e=8.23\times 10^{-8}\rm{N}$$

万有引力的大小为：

$$F_g = G \frac{m_1 m_2}{r^2}$$
$$G=6.67\times10^{-11} \rm{N}\cdot \rm{m}^2/\mathrm{kg}^2$$
$$m_1=m_p=1.67\times10^{-27}\rm{kg},m_2=m_e=9.11x10^{-31}\rm{kg}$$
$$  r=5.29\times10^{-11}\rm{m}$$
$$F_g=3.62\times10^{-47}\rm{N}$$

(2)

$$\frac{F_e}{F_g} \approx 2.27 \times 10^{39}$$

(3)

电子绕质子做圆周运动，向心力由库仑力提供(万有引力可忽略)：

$$\frac{m_e v^2}{r} = F_e = k \frac{e^2}{r^2}$$
$$v^2 = \frac{k e^2}{m_e r}$$
$$v=\sqrt{\frac{k e^2}{m_e r}}$$

$$v= \sqrt{\frac{8.99\times10^9\rm{N}\cdotp \rm{m}^2/ \rm{C}^2\times(1.60\times10^{-19}\rm{C})^2}{9.11x10^{-31}\rm{kg} \times 5.29\times10^{-11}\rm{m}}} \approx 2.19 \times 10^{6} \rm{m}/\rm{s}$$











\newpage
%1.1-5 卢瑟福实验证明：当两个原子核之间的距离小到 $10^{-15}$m 时,它们之间的排斥力仍遵守库仑定律. 
%金的原子核中有 79 个质子,氦的原子核(即 $\alpha$粒子)中有 2 个质子.已知每个质子带电 $e=1.60\times10^{-19}$ C,$\alpha$粒子的质量为$6.68\times10^{-27}$ kg.
%当$\alpha$粒子与金核相距为$6.9\times10^{-15}$ m 时(设这时它们都仍可当作点电荷),求：

%(1)$\alpha$粒子所受的力;

%(2)$\alpha$粒子的加速度.

1.1-5

(1)

库仑力的大小为：

$$F=k\frac{q_1q_2}{r^2}$$
$$k=8.99\times10^9\rm{N}\cdotp \rm{m}^2/ \rm{C}^2\qquad\qquad\qquad\qquad\qquad$$
$$ q_1=79e=79\times1.60\times10^{-19}\rm{C}=1.264\times10^{-17}$$
$$q_2=2e=2\times1.60\times10^{-19}\rm{C}=3.2\times10^{-17}\qquad $$
$$  r=6.9\times10^{-15}\rm{m}\qquad\qquad\qquad\qquad\qquad\qquad\quad$$
$$F\approx 764N$$

(2)


由牛顿第二定律有

$$F=ma$$

$$a = \frac{F}{m_\alpha} = \frac{764}{6.68 \times 10^{-27}}\approx 1.143 \times 10^{29} \rm{m}/\rm{s}$$


\newpage
%1.1-8 三个相同的点电荷放置在等边三角形的各顶点上.在此三角形的中心应放置怎样的电荷,才能使作用在每一点电荷上的合力为零？

1.1-8

设等边三角形的边长为 $a$,三个相同的点电荷 $q$ 分别位于顶点 $A, B, C$。三角形中心 $O$ 到每个顶点的距离为：

$$r = \frac{a}{\sqrt{3}}$$

考虑$A$所受的力

来自另外两个顶点电荷 $B$ 和 $C$ 的库仑力：

电荷 $B$ 对 $A$ 的力沿 $AB$ 方向,大小为：

$$F_{AB} = \frac{k q^2}{a^2}$$

同理 $C$ 对 $A$ 的力大小为：

$$F_{AC} = \frac{k q^2}{a^2}$$

这两个力夹角为 $60^\circ$,合力沿 $AO$ 的反方向(即指向三角形外),大小为：

$$F_A = 2 \cdot \frac{k q^2}{a^2} \cos 30^\circ = \frac{k q^2}{a^2} \cdot 2 \cdot \frac{\sqrt{3}}{2} = \frac{\sqrt{3} k q^2}{a^2}$$

方向沿 $OA$ 向外($O \to A$ 方向向外)。

考虑中心电荷 $Q$ 对顶点电荷 $A$ 的作用力

中心电荷 $Q$ 与顶点电荷 $q$ 的距离为 $r = a/\sqrt{3}$,库仑力大小为：

$$F_{O\to A} = \frac{k |Q q|}{r^2} = \frac{k |Q q|}{(a^2/3)} = \frac{3k |Q q|}{a^2}$$

方向沿 $OA$ 连线。若要使顶点电荷 $A$ 合力为零,则中心电荷对它的力必须与另外两个顶点电荷对它的力方向相反、大小相等。

为了平衡,中心电荷 $Q$ 对 $A$ 的力必须指向中心 $O$(即沿 $AO$ 方向从 $A$ 指向 $O$)。因此 $Q$ 与 $q$ 必须是异号,且力大小为：

$$\frac{3k |Q q|}{a^2} = \frac{\sqrt{3} k q^2}{a^2}$$
$$|Q| = \frac{\sqrt{3}}{3} q = \frac{q}{\sqrt{3}}$$

在三角形中心放置电荷 $-\frac{q}{\sqrt{3}}$,可使每个顶点电荷所受合力为零。




\newpage
%1.1-9 电荷量都是$Q$的两个点电荷相距为$l$,连线中点为$O$;有另一点电荷$q$,在连线的中垂面上距$O$为$x$处.
%(1)求$q$受的力；(2)若$q$开始时是静止的,然后让它自己运动,它将如何运动？分别就$q$与$Q$同号和异号两种情况加以讨论.

1.1-9

(1)

库仑力的大小为：

$$F=\frac{1}{4\pi\varepsilon_0}\frac{q_1q_2}{r^2}$$

$$PA = PB = \sqrt{\left(\frac{l}{2}\right)^2 + x^2} = \sqrt{\frac{l^2}{4} + x^2}$$

电荷 $ Q $ 在 $ A $ 对 $ q $ 的作用力大小：

$$F_A = \frac{1}{4\pi\varepsilon_0}\frac{q_1q_2}{r^2} = \frac{1}{4\pi\varepsilon_0}\frac{Q q}{\frac{l^2}{4} + x^2}$$

$$ \vec{F}_{A \to q} = \frac{1}{4\pi\varepsilon_0}\frac{Q q}{r^2} \cdot \frac{\overrightarrow{AP}}{|\overrightarrow{AP}|} \qquad \vec{F}_{B \to q} = \frac{1}{4\pi\varepsilon_0}\frac{Q q}{r^2} \cdot \frac{\overrightarrow{BP}}{|\overrightarrow{BP}|} $$

其中 $ r = \sqrt{(l/2)^2 + x^2} $

合力的水平分量：

来自 $ A $ 的分量：$ \frac{k Q q}{r^2} \cdot \frac{l/2}{r} $  

来自 $ B $ 的分量：$ \frac{k Q q}{r^2} \cdot \frac{l/2}{r} $  

两者抵消 ,合力水平分量为 0。

合力的中垂线分量：

来自 $ A $ 的分量：$ \frac{k Q q}{r^2} \cdot \frac{x}{r} $  

来自 $ B $ 的分量相同,因为 $ x $ 坐标一样。  

$$F = 2 \cdot \frac{1}{4\pi\varepsilon_0}\frac{ Q q}{r^2} \cdot \frac{x}{r} =\frac{Q q x}{2\pi\varepsilon_0r^3}$$

其中 $ r = \sqrt{\frac{l^2}{4} + x^2} $。

$$\vec{F} = \frac{Qqx}{2\pi\varepsilon_0(x^2+l^2/4)^{3/2}} \ \hat{y}$$

这里 $ \hat{y} $ 是中垂线方向(从 $ O $ 指向 $ P $ 的方向)

\newpage
(2) 

初始时 $ q $ 静止在 $ P $,设PO=x,$ x \neq 0 $。

情况 1：$ q $ 与 $ Q $ 同号($ Q q > 0 $)

由 (1) 知力为：

$$F =\frac{Qqx}{2\pi\varepsilon_0(x^2+l^2/4)^{3/2}} $$

当 $ Q q > 0 $ 时,$ q $ 受到一个沿中垂线远离 $ O $ 的力,

$$\frac{Qqx}{2\pi\varepsilon_0(x^2+l^2/4)^{3/2}} >0$$

$ q $ 将沿中垂线加速远离 $ O $ 点，直至无穷远,最终趋于匀速

情况 2：$ q $ 与 $ Q $ 异号($ Q q < 0 $)

此时力公式中 $ Q q < 0 $,所以：

$$\frac{Qqx}{2\pi\varepsilon_0(x^2+l^2/4)^{3/2}} $$

力的方向与$\vec{OP}$反向力指向 $ O $。,

初始 $ x_0 \neq 0 $,静止释放,受力指向 $ O $,所以 $ q $ 向 $ O $ 加速运动。  

到达 $ O $ 点时速度最大,$F=0$由于对称性,过 $ O $ 点后受力指向 $ O $(即减速)

对于大振幅,恢复力非线性,但仍为周期性振动,沿中垂线在 $ -x_0 $ 与 $ x_0 $ 之间振荡。

(i)若$q$与$Q$同号：$q$沿中垂线背离$O$运动至无穷远。

(ii)若$q$与$Q$异号：$q$以$O$为中心沿中垂线做往复振动。





\newpage
%1. 1-10 两小球质量都是 $m$,都用长为 $l$ 的细线挂在同一点；若它们带上相同的电荷量,平衡时两线夹角为 $2\theta$. 设小球的半径都可略去不计,求每个小球上的电荷量.

1.1-10

两球连线水平,每根线长 $ l $,与竖直方向夹角 $ \theta $,则：

$$\frac{d/2}{l} = \sin\theta$$
$$d = 2l\sin\theta$$

库仑力大小

$$F_e = \frac{1}{4\pi\varepsilon_0}\frac{q^2}{d^2} = \frac{1}{4\pi\varepsilon_0}\frac{q^2}{(2l\sin\theta)^2} = \frac{1}{4\pi\varepsilon_0}\frac{q^2}{4l^2\sin^2\theta}$$

对小球,水平方向：

$$T\sin\theta = F_e$$

竖直方向：

$$T\cos\theta = mg$$

两式相除：

$$\tan\theta = \frac{F_e}{mg}$$

代入 $ F_e $：

$$\tan\theta = \frac{q^2}{16\pi\varepsilon_0 l^2\sin^2\theta \cdot mg}$$


$$q^2 = 16\pi\varepsilon_0 l^2\sin^2\theta \cdot mg \tan\theta$$

$$q = \pm4l\sin\theta\sqrt{\pi\varepsilon_0mg\tan\theta}$$



\end{document}